\section{Tecnologías}\label{sec:tecnologias}

\subsection{Tecnologías identificadas}\label{subsec:tecnologias-identificadas}

Como resultado del proceso de búsqueda, selección y extracción de información se identificaron cuatro alternativas de software libre y de código abierto relacionadas con servicios de directorio.

Las tecnologías identificadas fueron:

\subsubsection{OpenLDAP}

OpenLDAP fue identificado como una de las alternativas de código abierto para la implementación de servicios de directorio. Según la documentación oficial del proyecto, se trata de una implementación de código abierto del protocolo \acs{LDAP}, orientada a entornos distribuidos y distribuida bajo la OpenLDAP Public License, una licencia libre de carácter permisivo con características similares a las licencias \ac{BSD} \citep{OpenLDAP2026}.

\subsubsection{389 Directory Server}

389 Directory Server fue identificado como una alternativa orientada a despliegues de mayor escala, incluyendo entornos corporativos y académicos. De acuerdo con la documentación oficial del proyecto, se trata de un servidor de directorio de nivel empresarial cuyo esquema de licenciamiento contempla una \acs{GPL} Exception y componentes distribuidos bajo Apache 2.0 y \acs{LGPL} para determinados elementos y extensiones \citep{DS3892026}.

\subsubsection{ApacheDS}

Apache Directory Server (ApacheDS) constituye una alternativa desarrollada en Java bajo el proyecto Directory de la Apache Software Foundation. Conforme a su documentación oficial, la solución se distribuye bajo Apache License 2.0 \citep{ApacheDS2026}.

\subsubsection{OpenDJ}

OpenDJ corresponde a una alternativa cuyos orígenes se remontan a los desarrollos de Sun Microsystems y ForgeRock, y que actualmente continúa su desarrollo por medio de la comunidad Open Identity Platform. La solución se distribuye bajo la \ac{CDDL} \citep{OpenDJ2026}.

\section{Incorporación de tecnología adicional}\label{sec:incorporacion-tecnologia-adicional}

Los resultados obtenidos mediante el \ac{SMS} permitieron establecer cuatro alternativas tecnológicas. Sin embargo, para ampliar el conjunto de soluciones consideradas en la etapa posterior de evaluación, se incorporó adicionalmente FreeIPA.

Es importante señalar que FreeIPA no constituye un resultado directo del Mapeo Sistemático de Literatura realizado en esta investigación. Por tanto, su incorporación no debe interpretarse como una quinta tecnología identificada mediante el proceso de búsqueda bibliográfica.

La incorporación de FreeIPA se realizó debido a que, de acuerdo con la documentación oficial del proyecto, se trata de una plataforma de código abierto orientada a la gestión centralizada de identidades y políticas, que integra diferentes componentes relacionados con los servicios de directorio y la autenticación \citep{FreeIPA2026}. Estas características hacen que constituya una alternativa tecnológica relevante para ser considerada junto con las soluciones obtenidas mediante el \ac{SMS}.

De esta manera, la incorporación de FreeIPA permite ampliar el conjunto de alternativas tecnológicas que serán sometidas posteriormente al proceso de evaluación, manteniendo diferenciada su procedencia respecto de las tecnologías identificadas directamente en el mapeo.

\section{Conjunto de tecnologías candidatas}\label{sec:conjunto-tecnologias-candidatas}

Como resultado de la identificación realizada mediante el \ac{SMS} y de la incorporación de la alternativa adicional, se estableció el siguiente conjunto de tecnologías candidatas, el cual se presenta en la \autoref{tab:tecnologias_candidatas}:

\begin{table}[htbp]
	\centering
	\caption{Conjunto de tecnologías candidatas.}
	\fontsize{9pt}{10pt}\selectfont
	\renewcommand{\arraystretch}{1.4} % Ajusta el espaciado vertical de las filas
	\begin{tabularx}{\textwidth}{|l|l|X|}
		\hline
		\rowcolor[gray]{0.85} % Sombreado gris del encabezado
		\multicolumn{1}{|c|}{\textbf{Tecnología}} &
		\multicolumn{1}{c|}{\textbf{Procedencia}} &
		\multicolumn{1}{c|}{\textbf{Licenciamiento}}
		\\
		\hline
		OpenLDAP                                  &
		Resultado del \acs{SMS}                   &
		OpenLDAP Public License
		\\
		\hline
		389 Directory Server                      &
		Resultado del \acs{SMS}                   &
		\ac{GPL} Exception / Apache 2.0 / \ac{LGPL}
		\\
		\hline
		ApacheDS                                  &
		Resultado del \acs{SMS}                   &
		Apache License 2.0
		\\
		\hline
		OpenDJ                                    &
		Resultado del \acs{SMS}                   &
		Common Development and Distribution License (\ac{CDDL})
		\\
		\hline
		FreeIPA                                   &
		Incorporación adicional                   &
		Código abierto
		\\
		\hline
	\end{tabularx}
	\label{tab:tecnologias_candidatas}
\end{table}
